% !TeX root = ../../main.tex
\subsection{can と be able to}
\begin{definitionbox}{\textbf{可能を表す助動詞}}
    \textbf{can} = \fitblankbf[]{be} \fitblankbf[]{able} \fitblankbf[]{to} + \fitblankbf[]{動詞の原形}

    \begin{flushright}
        「〜\fitblankbf[]{することができる}」
    \end{flushright}
\end{definitionbox}
\vspace{0.5em}

\subsection{can の過去形}
    can の過去形は\fitblankbf[]{could}である。また、will と can は似ているため、まとめて覚えよう。

    \begin{techniquebox}{\textbf{will と can の対応}}
        \begin{itemize}
            \item \makebox[28em][l]{\fitblankbf[]{will} = be going to + 動詞の原形}%
                $\Longrightarrow$ \fitblankbf[]{would}
            \item \makebox[28em][l]{\fitblankbf[]{can} = be able to + 動詞の原形}%
                $\Longrightarrow$ \fitblankbf[]{could}
        \end{itemize}
    \end{techniquebox}
    \vspace{0.5em}

    ※上記の be とは\fitblankbf[]{be 動詞}という意味であるので、主語の\fitblankbf[]{単複}や\fitblankbf[]{時制}によって\\
    \qquad \fitblankbf[]{適する形}に変化させなければならない。

    % \begin{enumerate}
    %     \item She can swim very fast.
    %     \dialogueblank{彼女はとても速く泳ぐことができる。}
    %     \vspace{0.5em}
    %     \item I am able to speak a little English.
    %     \dialogueblank{私は少し英語を話すことができる。}
    %     \vspace{0.5em}
    %     \item He could run ten kilometers when he was young.
    %     \dialogueblank{彼は若いころ10キロ走ることができた。}
    % \end{enumerate}

\subsection{助動詞 may}
    \settowidth{\dimen0}{\fitblankbf[]{May I} ～？}
    \makebox[\dimen0][l]{may }「〜 \fitblankbf[]{かもしれない} (\fitblankbf[]{推定})、
    \fitblankbf[]{してもよい} (\fitblankbf[]{許可})」
    \vspace{1em}

    May I ～？ 「～ \fitblankbf[]{していいですか} ？」

\subsection{may の過去形}
    may の過去形は\fitblankbf[]{might}である。

    % \begin{enumerate}
    %     \item It may rain this afternoon.
    %     \dialogueblank{今日の午後は雨が降るかもしれない。}
    %     \vspace{0.5em}
    %     \item May I use your pen?
    %     \dialogueblank{あなたのペンを使ってもいいですか。}
    % \end{enumerate}

\newpage
